%============================================================
% FORMATTING: BASE LAYER (SHARED)
% Aturan yang 100% sama di semua dokumen KP.
% Aturan per-dokumen (geometry, spacing) ada di formats/format_buku.tex
% yang di-input SETELAH file ini di main_buku.tex.
%============================================================

% -- Warna Institusi ------------------------------------------
\definecolor{pensorange}{RGB}{255, 149, 0}   % Orange KP (#FF9500)
\definecolor{penskuning}{RGB}{255, 193, 7} % Kuning/emas garis cover

% -- Nonaktifkan Hyphenation (pemenggalan kata) ----------------------
% Mencegah LaTeX memotong kata di akhir baris (mis. "Elek-tronika").
% tolerance=9999 + emergencystretch memberi ruang fleksibel agar
% LaTeX tidak perlu Overfull saat hyphenation dinonaktifkan.
\hyphenpenalty=10000
\exhyphenpenalty=10000
\tolerance=9999
\emergencystretch=3em

% -- Lebar Judul di Cover & Halaman Judul (harus sama) ---------
% 95mm = 0.92 * textwidth_normal (0.92 * 103mm = 94.76mm)
\newlength{\LebarJudulKP}
\setlength{\LebarJudulKP}{90mm}  % Silahkan diganti sesuai keinginan

% -- Indentasi & Paragraf -------------------------------------
% \parindent TIDAK diset di sini -- diset di formats/format_buku.tex
\setlength{\parskip}{0pt}

% -- Persamaan Matematika -------------------------------------
\setlength{\mathindent}{10mm}
\setlength{\abovedisplayskip}{12pt}
\setlength{\belowdisplayskip}{12pt}
\setlength{\abovedisplayshortskip}{12pt}
\setlength{\belowdisplayshortskip}{12pt}

% Nomor persamaan dicetak tebal-miring
\makeatletter
\renewcommand{\tagform@}[1]{\maketag@@@{\bfseries\itshape(#1)\@@italiccorr}}

% -- Macro: render fontsize tetap, cuma shrink kalau kepanjangan --
% Usage: \fitline{<max-width>}{<content>}
\newcommand{\fitline}[2]{%
  \setbox0=\hbox{{#2}}%
  \ifdim\wd0>#1
    \resizebox{#1}{!}{{#2}}%
  \else
    {#2}%
  \fi
}
\makeatother

% -- Format Judul Bab -----------------------------------------
% Nomor bab angka Arab: BAB 1, BAB 2, ... (Pasal 4.2.2 no. 3b)
\renewcommand{\thesection}{\arabic{chapter}.\arabic{section}}

% Mulai dari batas margin atas (spacing before = 0pt)
% Judul bab 1 spasi (12pt) di bawah tulisan BAB
% Kalimat awal mulai 2 spasi (24pt) di bawah judul bab
\titleformat{\chapter}[display]
{\centering\FontJudulBab\bfseries}
{\MakeUppercase{\chaptertitlename} \thechapter}
{0pt}
{\MakeUppercase}
\titlespacing*{\chapter}{0pt}{0pt}{12pt}

% -- Auto-selaraskan indentasi paragraf pertama dengan huruf judul (bukan nomor) --
\newsavebox{\KotakLabelJudul}
\newdimen\LebarLabelJudul
\newcommand{\SelaraskanIndentAsli}[1]{%
  \sbox{\KotakLabelJudul}{#1}%
  \global\LebarLabelJudul=\wd\KotakLabelJudul
  \global\advance\LebarLabelJudul by 1em
  \xdef\ParindentAsliTersimpan{\the\parindent}%
  \global\parindent=\LebarLabelJudul
  \global\everypar{\parindent=\ParindentAsliTersimpan\relax\global\everypar{}}%
}

% -- Format Judul Subbab (Section) ----------------------------
\titleformat{\section}
{\FontJudulSubbab\bfseries}
{\thesection}
{1em}
{}
[\SelaraskanIndentAsli{\thesection}]
\titlespacing*{\section}{0pt}{10pt}{3pt}

% -- Format Judul Sub-subbab (Subsection) ---------------------
\titleformat{\subsection}
{\FontJudulSubSubbab\bfseries}
{\thesubsection}
{1em}
{}
[\SelaraskanIndentAsli{\thesubsection}]
\titlespacing*{\subsection}{0pt}{8pt}{3pt}

% -- Format Judul Sub-sub-subbab (Subsubsection) --------------
% Tanpa nomor (secnumdepth=2) - heading ringan, label dikosongkan.
\titleformat{\subsubsection}
{\FontJudulSubSubSubbab\bfseries}
{}
{0pt}
{}
\titlespacing*{\subsubsection}{\parindent}{3pt}{2pt}

% -- Batas Penomoran dan Daftar Isi ---------------------------
% 0 = Chapter, 1 = Section, 2 = Subsection
\setcounter{secnumdepth}{2}
\setcounter{tocdepth}{2}

% -- Caption Gambar (Pasal 4.2.2 no. 6) -----------------------
% "Gambar 2.1: Judul gambar", label & judul bold (contoh pedoman
% hal. 26), di bawah gambar, menjorok 1 cm dari margin kiri-kanan,
% baris lanjutan sejajar judul. Sumber 8pt ditulis setelah caption.
\captionsetup[figure]{
  labelfont       = {bf},
  textfont        = {bf},
  font            = {normalsize},     % 10pt
  format          = hang,             % baris ke-2 sejajar judul (6c)
  justification   = raggedright,
  singlelinecheck = false,
  margin          = 1cm,              % menjorok 1 cm kiri-kanan (6b)
  skip            = 12pt,             % jarak 1 spasi (6i)
  name            = Gambar,
  labelsep        = colon             % "Gambar 2.1:"
}

% -- Caption Tabel (Pasal 4.2.2 no. 5) ------------------------
% "Tabel 2.1:" bold, judul 10pt normal, di atas tabel, menjorok
% 1 cm dari margin kiri-kanan, baris lanjutan sejajar judul.
\captionsetup[table]{
  labelfont       = {bf},             % "Tabel 2.1:" bold (5b)
  textfont        = {},               % judul 10pt normal (5e)
  font            = {normalsize},
  format          = hang,             % baris ke-2 sejajar judul (5c)
  justification   = raggedright,
  singlelinecheck = false,
  margin          = 1cm,              % menjorok 1 cm kiri-kanan (5b)
  skip            = 12pt,             % jarak 1 spasi (5i)
  position        = above,
  name            = Tabel,
  labelsep        = colon             % "Tabel 2.1:"
}
\renewcommand{\tablename}{Tabel}

% -- Format Default Tabularray --------------------------------
\SetTblrInner{
  hlines, vlines,
  row{1} = {bg=gray!20, font=\bfseries, halign=c},
  rowhead = 1,
}

\SetTblrInner[longtblr]{
  hlines, vlines,
  row{1} = {bg=gray!20, font=\bfseries, halign=c},
  rowhead = 1,
}

% -- Isi Tabel: lebih kecil (\small, +-9pt), spasi 1 ----------
% Berlaku otomatis untuk semua tblr/longtblr. Caption tidak terpengaruh.
\AddToHook{env/tblr/begin}{\small\linespread{1}\selectfont}
\AddToHook{env/longtblr/begin}{\small\linespread{1}\selectfont}

% -- Kustomisasi Khusus longtblr ------------------------------
% Caption longtblr disamakan dengan tabel biasa: "Tabel 2.1:" bold,
% judul 10pt normal, rata kiri.
\DefTblrTemplate{caption-sep}{default}{:\space}
\SetTblrStyle{caption-tag}{font=\normalsize\bfseries}
\SetTblrStyle{caption-text}{font=\normalsize}
\SetTblrStyle{caption-halign}{l}

\DefTblrTemplate{contfoot-text}{default}{}
\DefTblrTemplate{conthead-text}{default}{}
\DefTblrTemplate{capcont}{default}{}
\DefTblrTemplate{contfoot}{default}{}
\DefTblrTemplate{conthead}{default}{}

% -- Penomoran Gambar/Tabel/Persamaan per Bab -----------------
\renewcommand{\thefigure}{\arabic{chapter}.\arabic{figure}}
\renewcommand{\thetable}{\arabic{chapter}.\arabic{table}}
\renewcommand{\theequation}{\arabic{chapter}.\arabic{equation}}

% -- Header & Footer ------------------------------------------
\setlength{\headheight}{40pt}  % ruang untuk logo di header
\pagestyle{fancy}
\fancyhf{}
\fancyfoot[C]{\thepage}
\renewcommand{\headrulewidth}{0pt}


% Logo PENS untuk header Gaya 2 (diukur sekali agar lebar teks di
% sebelahnya bisa dihitung)
\newsavebox{\KotakLogoHeader}
\sbox{\KotakLogoHeader}{\includegraphics[height=1.0cm]{assets/pens/logo_pens.png}}

% Pagestyle kpruler dengan 2 macam gaya (Simple & PENS)
\fancypagestyle{kpruler}{
  \fancyhf{}
  \ifnum\StyleHeaderFooter=2
    % Gaya 2: PENS (Gambar 2)
    % Nama laporan & perusahaan dipindah ke header (baris ke-2) agar
    % nomor halaman bisa di bawah tengah (Pasal 4.2.2 no. 4c).
    % Logo di kiri, teks di ruang SISA sebelah kanan logo (rata tengah
    % di ruang itu) agar teks tidak menimpa logo. Baris yang terlalu
    % panjang otomatis dikecilkan (\fitline). Nama perusahaan di header memakai
    % \NamaPerusahaanSingkat (config/variables.tex) agar ringkas.
    \fancyhead[LO,LE]{%
      \raisebox{-0.2cm}{\usebox{\KotakLogoHeader}}\hspace{2mm}%
      \parbox[b]{\dimexpr\headwidth-\wd\KotakLogoHeader-2mm\relax}{%
        \centering\FontHeaderFooterPENS
        \fitline{\linewidth}{POLITEKNIK ELEKTRONIKA NEGERI SURABAYA}\\
        \fitline{\linewidth}{LAPORAN KERJA PRAKTIK~-~\MakeUppercase{\ifcsvoid{NamaPerusahaanSingkat}{\NamaPerusahaan}{\NamaPerusahaanSingkat}}}%
        \rule[-3mm]{0pt}{0pt}}}
    \renewcommand{\headrulewidth}{1.5pt}
    \fancyfoot[C]{\FontHeaderFooterPENS\thepage}
    \renewcommand{\footrulewidth}{1.5pt}
  \else
    % Gaya 1: Simple (Gambar 1)
    \fancyhead[C]{\FontHeaderFooterSimple\bfseries LAPORAN KERJA PRAKTIK\rule[-3mm]{0pt}{0pt}}
    \renewcommand{\headrulewidth}{1.5pt}
    \fancyfoot[C]{\FontHeaderFooterSimple\thepage}
    \renewcommand{\footrulewidth}{1.5pt}
  \fi
}

% Redefine plain to copy kpruler for consistency on chapter start pages
\makeatletter
\let\ps@plain\ps@kpruler
\makeatother

% -- Lingkungan Kutipan Panjang (lebih dari 2 baris) -----------
\newenvironment{kutipanpanjang}{%
  \begin{list}{}{%
      \setlength{\leftmargin}{10mm}
      \setlength{\rightmargin}{10mm}
      \setlength{\topsep}{6pt}
      \setlength{\itemsep}{0pt}
      \setlength{\parsep}{0pt}
    }
    \item\FontKutipanPanjang\ignorespaces
    }{%
    \unskip
  \end{list}
}

% -- Hyperref -------------------------------------------------
\hypersetup{
  colorlinks  = true,
  linkcolor   = black,
  citecolor   = black,
  urlcolor    = black,
  plainpages  = false,   % cegah duplicate anchor saat roman->arabic reset
  pdftitle    = {\JudulKP},
  pdfauthor   = {\NamaMahasiswaSatu}
}

% -- Nama Bahasa Indonesia ------------------------------------
\renewcommand{\contentsname}{DAFTAR ISI}
\renewcommand{\listfigurename}{DAFTAR GAMBAR}
\renewcommand{\listtablename}{DAFTAR TABEL}
\renewcommand{\bibname}{DAFTAR PUSTAKA}
\renewcommand{\appendixname}{LAMPIRAN}
\renewcommand{\chaptername}{BAB}

% -- Judul Daftar Isi/Gambar/Tabel: 12pt, Bold, Center --------
\renewcommand{\cfttoctitlefont}{\hfill\FontJudulDaftar\bfseries\MakeUppercase}
\renewcommand{\cftaftertoctitle}{\hfill\null}

\renewcommand{\cftloftitlefont}{\hfill\FontJudulDaftar\bfseries\MakeUppercase}
\renewcommand{\cftafterloftitle}{\hfill\null}

\renewcommand{\cftlottitlefont}{\hfill\FontJudulDaftar\bfseries\MakeUppercase}
\renewcommand{\cftafterlottitle}{\hfill\null}

\setlength{\cftbeforetoctitleskip}{0pt}
\setlength{\cftbeforeloftitleskip}{0pt}
\setlength{\cftbeforelottitleskip}{0pt}
\setlength{\cftaftertoctitleskip}{10pt}
\setlength{\cftafterloftitleskip}{10pt}
\setlength{\cftafterlottitleskip}{10pt}

% -- Prefix "Gambar" dan "Tabel" di LOF/LOT -------------------
\renewcommand{\cftfigpresnum}{Gambar }
\renewcommand{\cfttabpresnum}{Tabel }
% numwidth harus cukup lebar untuk menampung prefix terpanjang
% agar baris kedua (wrap) sejajar dengan teks, bukan kembali ke kiri.
% "Gambar 4.12 " ≈ 7em; "Tabel 4.13 " ≈ 6em
\setlength{\cftfignumwidth}{7.5em}
\setlength{\cfttabnumwidth}{6.5em}

% -- Titik-titik di Daftar Isi --------------------------------
\renewcommand{\cftdotsep}{1.5}

% -- Spasi Konsisten di Daftar Isi ----------------------------
\setlength{\cftbeforechapskip}{0pt}
\setlength{\cftbeforesecskip}{0pt}
\setlength{\cftbeforesubsecskip}{0pt}
\setlength{\cftbeforefigskip}{0pt}
\setlength{\cftbeforetabskip}{0pt}

% -- Font Level Bab: Bold -------------------------------------
\renewcommand{\cftchapfont}{\bfseries}
\renewcommand{\cftchappagefont}{\bfseries}
\renewcommand{\cftchapleader}{\cftdotfill{\cftdotsep}}

% -- Font Level Section: Normal -------------------------------
\renewcommand{\cftsecfont}{\normalfont}
\renewcommand{\cftsecpagefont}{\normalfont}

% -- Prefix "BAB" di Daftar Isi -------------------------------
\renewcommand{\cftchappresnum}{BAB }
% "BAB IV " ≈ 5em -- set numwidth eksplisit agar wrap sejajar teks
\setlength{\cftchapnumwidth}{5.5em}
% Indent section/subsection dikembalikan ke default agar menjorok bertingkat
% \setlength{\cftsecindent}{1.5em}
% \setlength{\cftsubsecindent}{3.8em}

% -- Jarak (Spasi) Atas/Bawah Gambar dan Tabel ----------------
\setlength{\intextsep}{10pt}
\setlength{\textfloatsep}{10pt}

\setlength{\bibsep}{3pt}

% -- Baris Yatim/Janda ----------------------------------------
% (Hyphenation sudah dinonaktifkan di bagian atas file ini.)
% Cegah satu baris paragraf tertinggal sendirian di bawah halaman
% (club/yatim) atau di atas halaman berikutnya (widow/janda).
\clubpenalty=10000
\widowpenalty=10000
\displaywidowpenalty=10000

% -- Halaman Kosong -------------------------------------------
% Halaman kosong pada cetak bolak-balik diberi keterangan
\makeatletter
\newif\if@mainmatter
\@mainmatterfalse
\def\cleardoublepage{%
  \clearpage
  \if@twoside
    \ifodd\c@page\else
      \hbox{}
      \vspace*{\fill}
      \begin{center}
        \textit{Halaman ini sengaja dikosongkan}
      \end{center}
      \vspace{\fill}
      \thispagestyle{kpruler}
      \newpage
      \if@twocolumn\hbox{}\newpage\fi
    \fi
  \fi
}

% -- Paksa Halaman Ganjil -------------------------------------
\newcommand{\cleartooddpage}{%
  \clearpage
  \ifodd\c@page\else
    \hbox{}
    \vspace*{\fill}
    \begin{center}
      \textit{Halaman ini sengaja dikosongkan}
    \end{center}
    \vspace{\fill}
    \thispagestyle{fancy}
    \newpage
  \fi
}
\makeatother
